\justifying
\NSFCBodyText

\subsubsection{已具备的实验条件}

依据申请人已确认的资源记录，可稳定使用 NVIDIA A40 共 8 张和 NVIDIA GeForce RTX 3090 Ti 共 6 张。现有 Linux、Python、深度学习框架、Git 和 LaTeX 工具链，可支持数据处理、区域模型训练、版本管理与结果整理。实验将按实际调度安排批次，研究预算不以所有显卡同时独占为前提。

数据方案以公开 CIESM 历史模拟、ERA5、ERA5-Land 和公开地形资料为主。当前具备按公开渠道获取这些资料并组织方法验证的条件；具体文件、变量频率与覆盖年份将在项目开始后逐项审计。现有事实材料未确认私有新能源场站数据或业务数值天气预报归档，本课题核心任务不依赖这些资源。

\subsubsection{尚缺少的实验条件及拟解决的途径}

尚需落实实验室合作人与最新 CIESM 数据及接口样例，即\textbf{【待补 F-CIESM-01：合作人、数据版本和联调安排】}。计划在第一季度对接数据格式、变量含义及实际应用需求，并在公开资料上先完成处理流程。内部数据的访问与使用按实验室要求执行，尚未落实前仅作为拟争取条件。

独立站点或第三方产品可用于增强真实性验证，但目前未获得本项目专属的授权证据。若能依法取得且与训练资料独立，则新增单独验证；在此之前，以资料内重建、跨区域迁移和一致性检查作为可执行研究范围，明确真实观测验证的不足。学院公共算力可作为补充申请资源，不预先计入固定配额。

\subsubsection{利用研究基地的计划与落实情况}

拟依托新疆大学现有科研环境开展算法研究，并与地球系统数值模拟教育部重点实验室围绕 CIESM 数据增强进行技术交流和应用对接。当前已确认的条件为申请人所在单位、现有计算资源及公开数据路径；实验室内部资源、合作人员、接口联调和成果正式接收仍按实际落实情况记录。项目将以可运行样例和应用证明支撑集成结论。
